\chapter{Introducción y motivación}\label{introduccion-y-motivacion}
\pagenumbering{arabic}

Este capítulo presenta el contexto del trabajo y la problemática que lo motiva. La sección~\ref{contexto-y-motivacion} introduce el contexto de aplicación y el objetivo del proyecto; las secciones~\ref{volumen-y-estructura} a~\ref{patrones-de-acceso} analizan las tres dimensiones de la problemática (volumen de datos, formato de almacenamiento y patrones de acceso); y la sección~\ref{oportunidad-sparse} expone la oportunidad técnica que explora el trabajo y la estructura del resto de la memoria.

\section{Contexto y motivación}\label{contexto-y-motivacion}

La necesidad de un almacenamiento eficiente para resultados de simulaciones hidráulicas surge de dos tendencias: el incremento de la resolución de los modelos hidrodinámicos y la demanda de análisis interactivos sobre volúmenes masivos de datos. Las simulaciones 2D de las ecuaciones de Saint-Venant \cite{saintvenant1871}, con formulaciones eficientes \cite{bates2010simple} e implementaciones sobre GPU \cite{morales2021triton}, permiten hoy modelar inundaciones a escala de cuenca con resolución de metros, generando series temporales de rásteres que alcanzan centenares de gigabytes por escenario. Los formatos habituales de salida no están diseñados para consultas analíticas sobre estos volúmenes. Plataformas como Google Earth Engine \cite{gorelick2017earthengine} han demostrado que desplazar el procesamiento al lugar donde residen los datos reduce drásticamente los tiempos de respuesta \cite{bauer2016earthsystem}, pero están orientadas a datos densos (imágenes satelitales) y carecen de soporte nativo para arrays dispersos, donde la mayor parte de las celdas están secas en cada instante.

El contexto de aplicación son las inundaciones fluviales. El simulador Triton \cite{triton} resuelve la propagación del agua sobre una malla regular y produce, para cada variable física y paso temporal, un fichero GeoTIFF (Geographic Tagged Image File Format) \cite{ogc_geotiff}: un formato ráster estándar en SIG (Sistemas de Información Geográfica) que integra los metadatos de georreferenciación en el propio fichero, pero diseñado para el acceso secuencial y no para el análisis de grandes series temporales. Los resultados permiten evaluar la extensión de la mancha de inundación, identificar zonas de riesgo y comparar escenarios. El \hyperref[anexo-c.-contexto-de-las-simulaciones-hidraulicas]{Anexo~C} amplía este contexto para lectores no familiarizados con la modelización hidráulica.

El objetivo principal de este TFG es diseñar e implementar un sistema de almacenamiento y análisis que permita almacenar estos resultados de forma compacta y eficiente, habilitando consultas analíticas rápidas y visualizaciones interactivas sin necesidad de cargar los rásteres completos en memoria. La Figura~\ref{fig:contexto-problema} resume el contraste entre el enfoque original y la solución propuesta: los GeoTIFF organizan los datos por variable y paso temporal (80 ficheros independientes), lo que obliga a lecturas masivas en serie para las consultas multi-temporales, mientras que el proceso ETL (\textit{Extract, Transform, Load}) los convierte en un array TileDB sparse 3D (tiempo, fila, columna) que almacena únicamente las celdas húmedas y responde las mismas consultas mediante una sola operación indexada.

\begin{figure}[ht]
\centering
\begin{tikzpicture}[
    node distance=6mm,
    arrow/.style={-Latex, thick, draw=black!60},
    etlarrow/.style={-Latex, line width=2pt, draw=blue!70!black},
]

%% —— Nodos principales (se definen antes de dibujar) ——
\node[rectangle, draw=gray!70, fill=black!5, rounded corners=3pt, thick,
      minimum width=50mm, minimum height=42mm] (geo) {};

\node[rectangle, draw=orange!80!black, fill=orange!8, rounded corners=3pt, thick,
      minimum width=50mm, minimum height=42mm, right=38mm of geo] (tdb) {};

%% —— Contenido bloque GeoTIFF ——
\node[font=\small\bfseries, above=2mm] at (geo.north) {GeoTIFF (por escenario)};
\foreach \r in {0,1,2,3} {
    \foreach \c in {0,1,2,3,4} {
        \node[rectangle, draw=gray!50, fill=white, minimum width=7mm, minimum height=5mm]
            at ([xshift=(\c-2)*8.5mm+3mm, yshift=(1.5-\r)*7.5mm] geo.center) {};
    }
}
% Etiquetas de variable: misma fórmula vertical que las filas de la rejilla
\foreach \r/\vn in {0/H, 1/QX, 2/QY, 3/MH} {
    \node[font=\tiny, anchor=east] at ([xshift=-19mm, yshift=(1.5-\r)*7.5mm] geo.center) {\vn};
}
\node[font=\tiny, below=3mm] at (geo.south) {4 variables $\times$ 20 pasos $=$ 80 ficheros $\approx$ 235 GB};

%% —— Contenido bloque TileDB ——
\node[font=\small\bfseries, above=2mm] at (tdb.north) {TileDB sparse};
\begin{scope}[shift={(tdb.center)}]
    \draw[fill=orange!15, draw=orange!70!black, thick] (-1.1cm,-0.75cm) rectangle (1.1cm,0.75cm);
    \draw[fill=orange!25, draw=orange!70!black, thick]
        (-1.1cm,0.75cm) -- (-0.75cm,1.1cm) -- (1.45cm,1.1cm) -- (1.1cm,0.75cm) -- cycle;
    \draw[fill=orange!20, draw=orange!70!black, thick]
        (1.1cm,0.75cm) -- (1.45cm,1.1cm) -- (1.45cm,-0.4cm) -- (1.1cm,-0.75cm) -- cycle;
    \node[font=\tiny\itshape] at (0, 0.15cm) {tiempo};
    \node[font=\tiny\itshape] at (0,-0.2cm) {fila $\times$ col};
\end{scope}
\node[font=\tiny, below=3mm] at (tdb.south) {Solo celdas húmedas $\approx$ 17{,}5 GB ($-$93\,\%)};

%% —— Cajas de consulta ——
\node[rectangle, draw=red!60, fill=red!5, rounded corners=2pt, dashed, thick,
      align=center, font=\scriptsize, minimum width=50mm, below=10mm of geo] (qgeo)
    {\textbf{Q7: hora de llegada del frente}\\
     Leer 20 ficheros en secuencia $\to$ varios minutos};

\node[rectangle, draw=green!60!black, fill=green!5, rounded corners=2pt, dashed, thick,
      align=center, font=\scriptsize, minimum width=50mm, below=10mm of tdb] (qtdb)
    {\textbf{Q7: misma consulta}\\
     1 operación sobre el array $\to$ 79{,}5 s};

%% —— Flechas ——
\draw[etlarrow] (geo.east) -- node[above, font=\small\bfseries, text=blue!70!black] {ETL}
                                node[below, font=\scriptsize, text=blue!60] {ingesta + compresión}
               (tdb.west);
\draw[arrow, draw=red!60]        (geo.south)  -- (qgeo.north);
\draw[arrow, draw=green!60!black] (tdb.south) -- (qtdb.north);

\end{tikzpicture}
\caption{Contexto del trabajo: del almacenamiento original en GeoTIFF al array TileDB sparse.}
\label{fig:contexto-problema}
\end{figure}
\FloatBarrier

La problemática de partida combina tres dimensiones que se refuerzan entre sí: un volumen de datos considerable, un formato de almacenamiento no diseñado para análisis y unos patrones de acceso que el formato no cubre con eficiencia.

\section{Volumen y estructura de los datos}\label{volumen-y-estructura}

Cada simulación con el simulador Triton produce 80 ficheros GeoTIFF georreferenciados (4 variables físicas $\times$ 20 pasos temporales), con un tamaño aproximado de 3 GB por fichero y un total cercano a 235 GB por escenario. El dominio espacial está discretizado en una malla regular de 34 200 $\times$ 23 040 píxeles, lo que supone aproximadamente 787 millones de celdas por paso temporal y unas 15 700 millones de celdas en la simulación completa. Procesar este volumen con herramientas convencionales exige cargar cada ráster completo en memoria, lo que requiere alrededor de 3 GB por variable y resulta inviable en estaciones de trabajo con menos de 32 GB de RAM.

\section{Inadecuación del formato GeoTIFF para consultas analíticas}\label{inadecuacion-geotiff}

GeoTIFF \cite{ogc_geotiff} está concebido para el acceso eficiente a un único ráster: no incorpora un motor de consultas, no indexa subconjuntos espaciales o temporales y no permite filtrar por valor sin descomprimir el ráster completo. Operaciones habituales como ``calcular la primera hora en que el agua llega a cada celda'' exigen recorrer y descomprimir los 20 pasos temporales: sobre 235 GB de datos crudos, esto supone tiempos del orden de minutos a horas y, en algunos casos, fallos por agotamiento de memoria.

\section{Patrones de acceso típicos del análisis hidráulico}\label{patrones-de-acceso}

El usuario final no consulta los datos como rásteres aislados, sino como una colección espacio-temporal coherente. Los patrones de acceso más habituales son cinco:

\begin{itemize}
\item \textbf{Consultas en un punto o zona}: Valor de calado en un punto y un instante concretos, o serie temporal completa en un punto.
\item \textbf{Consultas espaciales acotadas}: Estadísticos o mapas sobre una ventana geográfica reducida (por ejemplo, un núcleo urbano de 5$\times$5 km).
\item \textbf{Consultas multi-temporales}: Agregados sobre los 20 pasos temporales completos (hora de llegada del frente, duración total de la inundación, calado máximo por celda).
\item \textbf{Consultas de peligrosidad}: Clasificación de cada celda según criterios combinados de calado y caudal unitario (Russo et al. 2013 \cite{russo2013}).
\item \textbf{Consultas comparativas}: Contraste de métricas entre varios escenarios simulados (por ejemplo, año seco frente a año húmedo).
\end{itemize}

Tres de ellos (espaciales acotados, multi-temporales y peligrosidad) penalizan especialmente la lectura directa de GeoTIFF, y los cinco determinan el diseño de las 24 consultas del motor analítico.

\section{Oportunidad: dominio mayoritariamente seco}\label{oportunidad-sparse}

En una inundación típica solo el $\sim$8,3 \% de las celdas del dominio contiene agua en cada paso temporal ($\sim$65 millones de las $\sim$787 millones totales); el resto son celdas secas sin información útil. Esta característica convierte el almacenamiento disperso (que solo guarda las celdas no nulas y sus coordenadas) en un candidato natural, y constituye la oportunidad técnica que explora este trabajo.

Tras este capítulo introductorio, el capítulo~\ref{descripcion-de-los-datos} describe los datos de entrada; el capítulo~\ref{analisis-de-tecnologias} analiza las tecnologías evaluadas y el funcionamiento de TileDB; el capítulo~\ref{diseno-del-sistema} presenta el diseño del sistema, el proceso ETL y la elección experimental de la configuración de almacenamiento; el capítulo~\ref{herramientas-desarrolladas} las herramientas desarrolladas (motor de consultas, comparativa frente a GeoTIFF y aplicación web); el capítulo~\ref{validacion-del-sistema} la validación y el rendimiento; y el capítulo~\ref{conclusiones-y-trabajo-futuro} las conclusiones. Los Anexos A--E completan la memoria con el catálogo de consultas, las tablas detalladas del \textit{benchmark}, el contexto hidráulico, el proceso ETL detallado y la guía de despliegue.

